% Editable TikZ source for the 8-bit parity checker.
% The seven XOR operations each use four NAND2 cells in the implementation.
\begin{tikzpicture}[x=1.12cm,y=.63cm,
  xorcell/.style={draw=reportblue,fill=blue!4,minimum width=1.15cm,minimum height=.58cm,font=\small},
  wire/.style={-{Latex[length=2mm]},line width=.55pt},font=\small]
\node[xorcell] (x0) at (3,7) {XOR};
\node[xorcell] (x1) at (3,5) {XOR};
\node[xorcell] (x2) at (3,3) {XOR};
\node[xorcell] (x3) at (3,1) {XOR};
\node[xorcell] (x4) at (6.3,6) {XOR};
\node[xorcell] (x5) at (6.3,2) {XOR};
\node[xorcell] (x6) at (9.6,4) {XOR};
\foreach \i/\target/\yy in {0/x0/7.45,1/x0/6.55,2/x1/5.45,3/x1/4.55,4/x2/3.45,5/x2/2.55,6/x3/1.45,7/x3/.55}{
 \node[anchor=east] at (.4,\yy) {$A_{\i}$};
 \draw[wire] (.55,\yy)--(\target.west);
}
\draw[wire] (x0.east)--(x4.west);
\draw[wire] (x1.east)--(x4.west);
\draw[wire] (x2.east)--(x5.west);
\draw[wire] (x3.east)--(x5.west);
\draw[wire] (x4.east)--(x6.west);
\draw[wire] (x5.east)--(x6.west);
\draw[wire] (x6.east)--(11,4) node[right] {$Y$};
\node[anchor=north,font=\footnotesize,align=center] at (5.8,-.1)
 {Each XOR = 4 NAND2 cells\qquad Total: 28 NAND2 cells, 112 MOSFETs};
\end{tikzpicture}
